\title{DeepSeekMath: Pushing the Limits of Mathematical Reasoning in Open Language Models}

\begin{abstract}
The intricate and highly structured requirements of mathematical reasoning remain a formidable obstacle for language models. In this study, we present DeepSeekMath~7B, a model developed by further pre-training DeepSeek-Coder-Base-v1.5 7B with a dataset containing 120B math-focused tokens extracted from Common Crawl, along with additional natural language and programming code sources. DeepSeekMath~7B attains a remarkable 51.7\% accuracy on the challenging, competition-oriented MATH benchmark, accomplishing this without dependence on auxiliary toolkits or ensemble voting, thereby nearing the results achieved by Gemini-Ultra and GPT-4. When applying self-consistency across 64 generated solutions, DeepSeekMath~7B reaches 60.9\% on the MATH evaluation. The model’s advancements in mathematical reasoning are rooted in two primary innovations: firstly, the implementation of a comprehensive data selection framework that systematically leverages abundant web-based resources; and secondly, the development of Group Relative Policy Optimization (GRPO)—an enhanced form of Proximal Policy Optimization (PPO)—which both boosts the model’s capacity for mathematical inference and improves memory efficiency in PPO-based training.
\end{abstract}

\section{Introduction}

Large language models (LLMs) have dramatically transformed the landscape of mathematical reasoning within artificial intelligence, driving remarkable progress on tasks such as the quantitative reasoning benchmark \citep{MATH} and the geometry reasoning benchmark \citep{trinh2024solving}. In addition, LLMs have demonstrated valuable utility in supporting humans as they tackle intricate mathematical challenges \citep{tao}. Nonetheless, the most advanced models to date, including GPT-4 \citep{gpt4} and Gemini-Ultra \citep{gemini}, are not freely accessible to the public, while existing open-source alternatives exhibit a marked gap in performance relative to their proprietary counterparts.

This paper presents DeepSeekMath, a specialized language model tailored for mathematical reasoning, which achieves substantially better performance than existing open-source alternatives and nears the proficiency of GPT-4 on standardized academic tasks. Central to this advancement is the construction of the DeepSeekMath Corpus: an extensive, high-quality pre-training dataset consisting of 120B math-related tokens. The corpus is curated from Common Crawl (CC) data through the use of a classifier built on the fastText architecture \citep{joulin2016fasttext}. During its initial development, the classifier leverages samples from OpenWebMath \citep{openwebmath} as positive examples, while integrating a variety of other web data as negative samples. The model is then applied to extract further positive samples from CC, which are subsequently subjected to human annotation for greater accuracy. The classifier undergoes retraining with this improved dataset, leading to enhanced classification results. Evaluation outcomes demonstrate that the resultant large-scale corpus maintains high quality: the DeepSeekMath-Base 7B model secures a score of 64.2\% on GSM8K \citep{gsm8k} and 36.2\% on the MATH competition benchmark \citep{MATH}, surpassing Minerva 540B \citep{minerva}. Furthermore, due to the multilingual nature of the DeepSeekMath Corpus, noticeable gains are observed in Chinese mathematical benchmarks \citep{wei2023cmath,agieval}. We suggest that our methodology for mathematical data collection offers a valuable foundation for further research and that there exists considerable potential for subsequent refinements.

DeepSeekMath-Base builds upon DeepSeek-Coder-Base-v1.5 7B \citep{deepseek-coder}, since initializing from a model trained on code yields superior results relative to using a generic LLM. Additionally, our findings indicate that mathematical training boosts performance not only in mathematics-related domains but also in benchmarks like MMLU \citep{mmlu} and BBH \citep{bbh}. This suggests the approach enhances both mathematical proficiency and overall reasoning skills.

Following the pre-training phase, we perform mathematical instruction tuning on DeepSeekMath-Base utilizing datasets featuring chain-of-thought \citep{cot}, program-of-thought \citep{pot,pal}, and tool-integrated reasoning \citep{tora} methodologies. The tuned model, DeepSeekMath-Instruct 7B, outperforms all other 7B parameter models and demonstrates performance on par with instruction-tuned open-source models possessing 70B parameters.

In addition, we present Group Relative Policy Optimization (GRPO), a reinforcement learning (RL) algorithm inspired by Proximal Policy Optimization (PPO) \citep{schulman2017proximal}, which differs by omitting the critic network and instead leveraging group-based score baselines for enhanced efficiency. This modification results in a notable decrease in computational demands during training. Utilizing only a fraction of English instruction tuning data, GRPO achieves marked performance gains over the robust DeepSeekMath-Instruct baseline—demonstrated by improvements in both in-domain benchmarks (GSM8K: 82.9\% $\rightarrow$ 88.2\%, MATH: 46.8\% $\rightarrow$ 51.7\%) and out-of-domain tasks such as CMATH (84.6\% $\rightarrow$ 88.8\%) during reinforcement learning. We also propose a cohesive framework that unifies various approaches, including Rejection Sampling Fine-Tuning (RFT) \citep{yuan2023scaling}, Direct Preference Optimization (DPO) \citep{dpo}, PPO, and GRPO, highlighting that these approaches are fundamentally variants of either direct or simplified RL strategies. To further illuminate the core aspects of this framework, we carry out comprehensive empirical studies, investigating dimensions such as online versus offline optimization, outcome versus process supervision, as well as single-turn and iterative RL. Finally, we clarify the mechanisms by which our RL approach enhances the effectiveness of instruction-tuned models, and outline future avenues for advancing RL through the lens of this unifying paradigm.

\subsection{Contributions}
We offer advancements in scalable math pre-training and investigate the application and effects of reinforcement learning within this context.

\textbf{Math Pre-Training at Scale}

\begin{itemize}[topsep=0pt]
    \item Our study demonstrates that the publicly available Common Crawl dataset holds significant potential for mathematical applications. Through a carefully structured data selection process, we develop the DeepSeekMath Corpus—a curated, high-quality dataset encompassing 120B tokens extracted from web sources identified for their mathematical content. This corpus substantially surpasses previous efforts, being nearly 7 times larger than the math web pages included in Minerva \citep{minerva} and approximately 9 times bigger than the newly introduced OpenWebMath \citep{openwebmath}.

\item DeepSeekMath-Base 7B, our pre-trained base model, demonstrates performance on par with Minerva 540B \citep{minerva}. This suggests that model size alone does not solely determine mathematical reasoning abilities. Robust outcomes can also be obtained from smaller models when they are pre-trained with high-quality datasets.

\item We present the results obtained from our experiments on math training. Pretraining models on code enhances their performance in solving mathematical tasks, regardless of whether external tools are involved. This provides a partial response to the ongoing debate: \textit{does code training enhance reasoning skills?} Our evidence suggests that it does, at least within the domain of mathematical reasoning.

\item While the use of arXiv papers as training material is prevalent, particularly among mathematics-focused research, this approach does not yield significant gains across any of the mathematical benchmarks considered in this study.
\end{itemize}

\textbf{Exploration and Analysis of Reinforcement Learning}

\begin{itemize}[topsep=0pt]
    \item We present Group Relative Policy Optimization (GRPO), a novel and efficient reinforcement learning method. Unlike traditional approaches such as Proximal Policy Optimization (PPO), GRPO eliminates the use of a critic network and instead derives the baseline using group-based score estimation, thereby markedly lowering training resource consumption.
    \item Through empirical studies, we show that applying GRPO leads to substantial improvements in the capabilities of our instruction-tuned system, DeepSeekMath-Instruct, relying exclusively on the instruction-tuning corpus. Additionally, we identify improvements in generalization to out-of-domain data throughout the reinforcement learning phase.
    \item We establish a comprehensive framework for interpreting and connecting different algorithms, including RFT, DPO, PPO, and GRPO. A broad range of experimental analyses are conducted—such as comparing online and offline paradigms, evaluating supervision at the outcome versus process levels, and contrasting single-step with iterative reinforcement learning—to thoroughly explore the fundamental components of this overarching framework.
    \item Leveraging our consolidated framework, we delve into the mechanisms underpinning the success of reinforcement learning in large language models (LLMs) and propose several promising avenues for advancing the effectiveness of LLM reinforcement learning.
\end{itemize}

\subsection{Summary of Evaluations and Metrics}

\begin{itemize}[topsep=0pt]
    \item \textbf{English and Chinese Mathematical Reasoning}: We systematically evaluate our models across a variety of English and Chinese benchmark datasets, encompassing mathematical tasks that range from primary school to university difficulty. The English set comprises GSM8K \citep{gsm8k}, MATH \citep{MATH}, SAT \citep{llemma}, OCW Courses \citep{minerva}, and MMLU-STEM \citep{mmlu}. For Chinese, the benchmarks include MGSM-zh \citep{mgsm}, CMATH \citep{wei2023cmath}, Gaokao-MathCloze \citep{agieval}, and Gaokao-MathQA \citep{agieval}. Our evaluation focuses both on the model’s proficiency at producing complete, standalone textual explanations without external tools, and its capacity to tackle problems programmatically using Python.

On English evaluation tasks, DeepSeekMath-Base matches the performance of the proprietary Minerva 540B model \citep{minerva}, while outperforming all existing open-source base models—such as Mistral 7B \citep{mistral} and Llemma-34B \citep{llemma}—whether or not those models have undergone additional mathematical pre-training, frequently by a large margin. Particularly, DeepSeekMath-Base exhibits clear advantages on Chinese benchmarks, which can be attributed to our deliberate choice not to restrict math pre-training data to English as in prior studies \citep{minerva,llemma}, but rather to include high-quality datasets in multiple languages. Upon applying mathematical instruction tuning and reinforcement learning, our enhanced models, DeepSeekMath-Instruct and DeepSeekMath-RL, display remarkable capabilities, achieving for the first time in the open-source domain an accuracy exceeding 50\% on the competition-level MATH dataset.

\item \textbf{Formal Mathematics}: We assess DeepSeekMath-Base's ability on the informal-to-formal theorem proving task as described in \citep{dsp_proof}, utilizing the miniF2F benchmark \citep{minif2f} with the Isabelle proof assistant \citep{isabelle}. The results indicate that DeepSeekMath-Base achieves robust performance in few-shot autoformalization settings.
\item \textbf{Natural Language Understanding, Reasoning, and Code}: To obtain a well-rounded evaluation of general comprehension, reasoning, and coding abilities, we test DeepSeekMath-Base across several benchmarks: the Massive Multitask Language Understanding (MMLU) dataset \citep{mmlu} featuring 57 multiple-choice tasks spanning numerous disciplines; BIG-Bench Hard (BBH) \citep{bbh}, a suite of 23 complex problems predominantly involving multi-step reasoning; as well as HumanEval \citep{codex} and MBPP \citep{mbpp}, two widely recognized standards for code generation models. Pre-training on mathematical data leads to enhanced performance in both language understanding and reasoning tasks.

\section{Math Pre-Training}

\subsection{Data Collection and Decontamination} This section details the methodology employed to assemble the DeepSeekMath Corpus utilizing data from Common Crawl. As illustrated in Figure \ref{fig:data_collect}, we introduce an iterative pipeline that enables the structured extraction of an extensive mathematical dataset, beginning with an initial seed corpus comprising a limited yet high-quality set of math-related resources. Importantly, this methodology is transferable and can be adapted to other areas, such as code datasets.

We begin by selecting OpenWebMath \citep{openwebmath}, an extensive repository of high-caliber mathematical web documents, to serve as our seed dataset. With this dataset, we proceed to train a fastText model \citep{joulin2016fasttext} in order to retrieve additional mathematical web content exhibiting similarities to OpenWebMath. For model training, 500,000 samples are randomly drawn from the seed as positive instances, while an equal number of web pages from Common Crawl are utilized as negative samples. Training is carried out using an open-source implementation\footnote{\url{https://fasttext.cc}}, with hyperparameters set as follows: vector dimensionality of 256, a learning rate of 0.1, a maximum word n-gram size of 3, a minimum word occurrence threshold of 3, and 3 epochs. To curtail the vastness of Common Crawl, we leverage both URL deduplication and near-duplicate filtering, which reduces the dataset to 40B HTML web pages. The trained fastText model is subsequently used to identify mathematical web pages from this deduplicated corpus. In an effort to maintain only high-quality mathematical data, the retrieved pages are ranked according to their fastText-generated relevance scores, retaining exclusively the highest-scoring documents. The resulting data volume is evaluated via pre-training experiments on token subsets of 40B, 80B, 120B, and 160B. In our initial pass, we opt to preserve only the top 40B tokens.

Following the initial round of data acquisition, a considerable number of mathematical web pages remain uncovered, largely due to the limited heterogeneity in the positive sample set used for fastText model training. To address this, we expand the seed corpus by incorporating additional sources of mathematical content. This begins with partitioning the full Common Crawl dataset by domain, where each domain encompasses web pages under a shared base URL. For every domain, we assess the proportion of web pages harvested during the initial collection. Domains for which more than 10\% of their pages have already been gathered are designated as mathematics-related (e.g., \url{mathoverflow.net}).

Following this, we engage in a manual review and annotation of the URLs containing mathematical material within the identified domains (for example, \url{mathoverflow.net/questions}). Any remaining uncollected web pages corresponding to these URLs are appended to the seed corpus. This strategy allows us to augment the number of positive examples, supporting the development of a more robust fastText model with improved capability to identify mathematical content in the next round. After completing four iterations, the process yields a total of 35.5M mathematical web pages, accumulating to 120B tokens. During the fourth iteration, we observe that nearly 98\% of all data was already gathered in the prior iteration; consequently, data collection is discontinued at this stage.

In order to prevent benchmark contamination, we adopt the approach described in \cite{deepseek-coder}, excluding any web content that includes questions or answers sourced from English mathematical benchmarks like GSM8K \citep{gsm8k} and MATH \citep{MATH}, as well as Chinese benchmarks such as CMATH \citep{wei2023cmath} and AGIEval \citep{agieval}. Specifically, our filtering process removes any portion of text that contains a 10-gram sequence that perfectly coincides with a sub-string found in these benchmarks. For benchmark passages that are under 10 grams but at least 3 grams in length, we apply precise string matching to identify and eliminate any compromised web pages from the math training dataset.

\subsection{Validating the Quality of the DeepSeekMath~Corpus}

We run pre-training experiments to investigate how the DeepSeekMath~Corpus is compared with the recently released math-training corpora:

\begin{itemize}[topsep=0pt]
    \item \textbf{MathPile} \citep{mathpile}: a comprehensive dataset encompassing 8.9B tokens, assembled from sources such as textbooks, Wikipedia, ProofWiki, CommonCrawl, StackExchange, and arXiv, with more than 85\% of the material originating from arXiv;
    \item \textbf{OpenWebMath} \citep{openwebmath}: a subset of CommonCrawl specifically filtered for mathematical text, comprising 13.6B tokens;
    \item \textbf{Proof-Pile-2} \citep{llemma}: an extensive mathematical collection made up of OpenWebMath, AlgebraicStack (which contains 10.3B tokens of math-related code), and a set of arXiv articles (28.0B tokens). For experiments utilizing Proof-Pile-2, we adopt the arXiv:Web:Code distribution of 2:4:1, consistent with the methodology described in \cite{llemma}.
\end{itemize}

\subsubsection{Training Setting}
\label{sec:quality-policy}
We perform mathematical instruction tuning on a base language model comprised of 1.3B parameters, architecturally consistent with the DeepSeek LLMs \citep{deepseek-llm}, hence referred to as DeepSeek-LLM 1.3B. For each distinct mathematical dataset, the model undergoes training for a total of 150B tokens. All training runs utilize the HAI-LLM \citep{haillm} framework, selected for its efficiency and lightweight design. Consistent with established DeepSeek LLM training protocols, we adopt the AdamW optimizer \citep{adamW} and configure its hyperparameters to $\beta_1=0.9$, $\beta_2=0.95$, and $\mathrm{weight\_decay}=0.1$. The learning rate is managed using a multi-stage schedule: it is linearly increased to its peak over the first 2,000 warmup steps, decays to 31.6\% of its maximum by the 80\% mark of training completion, and is further reduced to 10.0\% of its peak value upon reaching 90\% of total training. The upper bound for the learning rate is set at 5.3e-4. Training is performed with a token batch size of 4 million and a maximum context window of 4,000 tokens.

\subsubsection{Evaluation Results}

\textbf{The DeepSeekMath~Corpus is of high quality, covers multilingual mathematical content, and is the largest in size.}

\begin{itemize}[topsep=0pt]
    \item \textbf{High-quality}:
    We assess the downstream capabilities of our models across 8 mathematical benchmarks utilizing few-shot chain-of-thought prompting as outlined in \cite{cot}. Table \ref{tab:corpora_comparison} illustrates that models trained using the DeepSeekMath~Corpus significantly outperform others. Furthermore, as visualized in Figure \ref{fig:corpora_comparisons}, models leveraging the DeepSeekMath Corpus excel over those trained on Proof-Pile-2 at the same scale of 50B tokens (corresponding to a full epoch of Proof-Pile-2). This highlights the superior average quality of data in the DeepSeekMath Corpus.
    \item \textbf{Multilingual}:
    The DeepSeekMath~Corpus contains data in various languages, with English and Chinese being the most prevalent. Table \ref{tab:corpora_comparison} demonstrates that utilizing the DeepSeekMath~Corpus in training consistently leads to enhanced mathematical reasoning abilities in both English and Chinese. Conversely, prior mathematical corpora, which overwhelmingly focus on English, yield minimal gains or can even impair Chinese-language mathematical reasoning.
    \item \textbf{Large-scale}:
    DeepSeekMath~Corpus notably surpasses previous mathematical corpora in size. Figure \ref{fig:corpora_comparisons} shows that training DeepSeek-LLM 1.3B on this corpus results in a more accelerated and sustained performance improvement. In comparison, conventional corpora are much smaller and are reused repeatedly throughout training, causing model improvements to plateau rapidly.
\end{itemize}

\subsection{Training and Evaluating DeepSeekMath-Base 7B}

This section presents DeepSeekMath-Base 7B, a foundational model demonstrating notable proficiency in mathematical reasoning. We initialized this model using DeepSeek-Coder-Base-v1.5 7B \citep{deepseek-coder} and proceeded to train it over 500B tokens. The training dataset consists of 56\% DeepSeekMath Corpus, 4\% AlgebraicStack, 10\% arXiv sources, 20\% code from Github, and 10\% natural language content from Common Crawl, equally covering English and Chinese. Training configurations primarily follow the setup detailed in Section \ref{sec:quality-policy}, with the exception that the learning rate's peak is set at 4.2e-4 and the batch size is adjusted to 10M tokens.

We present an in-depth evaluation of DeepSeekMath-Base 7B’s mathematical proficiency, emphasizing its performance in generating independent mathematical solutions unaided by external resources, its effectiveness in addressing mathematical questions with tool assistance, and its competence in formal theorem verification. In addition to the mathematical focus, we offer a broader characterization of the base model, assessing its capabilities in natural language comprehension, reasoning abilities, and programming expertise.

\paragraph{Mathematical Problem Solving with Step-by-Step Reasoning}

We assess DeepSeekMath-Base's ability to solve mathematical problems through few-shot chain-of-thought prompting \citep{cot} on eight distinct benchmarks in both English and Chinese. The evaluation set spans tasks in quantitative reasoning—such as GSM8K \citep{gsm8k}, MATH \citep{MATH}, and CMATH \citep{wei2023cmath}—as well as multiple-choice assessments, including MMLU-STEM \citep{mmlu} and Gaokao-MathQA \citep{agieval}. Collectively, these benchmarks cover a wide range of mathematical disciplines and difficulty levels, from elementary school topics to those encountered in undergraduate studies.

Table \ref{tab:base_math_cot} demonstrates that among all open-source base models evaluated—including the prominent Mistral 7B \citep{mistral} and the newly introduced Llemma 34B \citep{llemma}, both of which have been trained on Proof-Pile-2 \citep{llemma}—DeepSeekMath-Base 7B consistently achieves the highest scores across eight benchmarks. Significantly, DeepSeekMath-Base attains more than a 10\% absolute improvement over prior open-source base models on the challenging competition-grade MATH dataset. Furthermore, it exceeds the performance of Minerva 540B \citep{minerva}, a proprietary base model with 77 times more parameters, which is constructed atop PaLM \citep{palm} and subsequently fine-tuned on mathematical corpora.

\paragraph{Mathematical Problem Solving with Tool Use} We assess model-assisted mathematical reasoning abilities on the GSM8K and MATH benchmarks using few-shot program-of-thought prompting \citep{pot,pal}. For each task, models are instructed to generate a Python script that may employ libraries such as \textit{math} and \textit{sympy} to handle complex calculations. The program’s output is treated as the predicted solution. Table \ref{tab:base_math_pot_proof} demonstrates that DeepSeekMath-Base 7B achieves superior results compared to the previous best-performing Llemma 34B model.

\paragraph{Formal Mathematics}

Automating formal proofs offers significant advantages in terms of improving the correctness, dependability, and productivity of mathematical verification, and has recently received growing interest. In our study, we assess DeepSeekMath-Base 7B on the informal-to-formal proof generation task introduced in \citep{dsp_proof}, which entails producing a formal proof given an informal problem statement, its formal equivalent, and an accompanying informal proof. For our experiments, we use the miniF2F benchmark \citep{minif2f}, which focuses on formalized versions of high-level Olympiad mathematics problems, prompting the model with a few examples to create Isabelle proofs for each instance. Adhering to the methodology of \cite{dsp_proof}, we have models draft proof sketches and subsequently rely on the established automated prover Sledgehammer \citep{sledgehammer} to address incomplete proof segments. Table \ref{tab:base_math_pot_proof} presents evidence that DeepSeekMath-Base 7B achieves noteworthy results in the task of proof autoformalization.

\paragraph{Natural Language Understanding, Reasoning, and Code}

We assess the capabilities of the model in natural language understanding using MMLU \citep{mmlu}, in reasoning using BBH \citep{bbh}, and in programming via HumanEval \citep{codex} and MBPP \citep{mbpp}. According to Table \ref{tab:understanding_reasoning_code}, DeepSeekMath-Base 7B demonstrates marked improvement over its predecessor, DeepSeek-Coder-Base-v1.5 \citep{deepseek-coder}, specifically in the domains of MMLU and BBH, underscoring the beneficial influence of math-oriented training on both language comprehension and reasoning proficiency. Moreover, through the integration of code tokens in ongoing training, DeepSeekMath-Base 7B is able to preserve the coding benchmark performance established by DeepSeek-Coder-Base-v1.5. In aggregate, DeepSeekMath-Base 7B achieves substantially better results than the general-purpose Mistral 7B \citep{mistral} across the trio of reasoning and coding evaluation tasks.

\section{Supervised Fine-Tuning}

\subsection{SFT Data Curation}

We develop a comprehensive mathematical instruction-tuning dataset that spans both English and Chinese problems across a range of mathematical domains and difficulty levels. Each problem is accompanied by solutions articulated in chain-of-thought (CoT) \citep{cot}, program-of-thought (PoT) \citep{pot,pal}, and tool-integrated reasoning formats \citep{tora}. In total, the dataset comprises 776K training examples.

\begin{itemize}[topsep=0pt]
    \item \textbf{English mathematical datasets}: Our annotation efforts extend to GSM8K and MATH, for which we provide tool-assisted solution annotations. Additionally, we include a selected portion of MathInstruct \citep{MathInstruct} and utilize the training split of Lila-OOD \citep{lila}, where each problem is addressed using either CoT or PoT strategies. The assembled English datasets represent a broad range of mathematical areas, such as algebra, probability, number theory, calculus, and geometry.
    \item \textbf{Chinese mathematical datasets}: We assemble a dataset of Chinese K-12 math questions encompassing 76 different sub-domains, including topics like linear equations. Each problem is annotated with both CoT and tool-integrated solution formats.
\end{itemize}

\subsection{Training and Evaluating DeepSeekMath-Instruct 7B}

In this section, we present DeepSeekMath-Instruct 7B, a model subjected to mathematical instruction tuning leveraging the DeepSeekMath-Base architecture. During training, examples are combined at random up to a total of 4,000 tokens per input. The model is then trained over 500 iterations, utilizing a batch size of 256 and maintaining a fixed learning rate of 5e-5 throughout the process.

We evaluate models' mathematical performance both without and with tool use, on 4 quantitative reasoning benchmarks in English and Chinese. We benchmark our model against the leading models of the time:

\begin{itemize}[topsep=0pt]
    \item \textbf{Closed-source models} comprise: (1) the GPT series, with GPT-4 \citep{gpt4} and GPT-4 Code Interpreter~\footnote{\url{https://openai.com/blog/chatgpt-plugins\#\#code-interpreter}} currently representing the most advanced, (2) Gemini Ultra and Pro \citep{gemini}, (3) Inflection-2 \citep{inflection-2}, (4) Grok-1~\footnote{\url{https://x.ai/model-card}}, along with several newer models from Chinese developers such as (5) Baichuan-3~\footnote{\url{https://www.baichuan-ai.com}}, (6) and the latest GLM-4~\footnote{\url{https://open.bigmodel.cn/dev/api\#glm-4}}.

    from the GLM family \citep{glm}.

These models are designed for broad applications, and the majority have been subjected to various alignment strategies. 
    \item \textbf{Open-source models} encompass: general-purpose large language models, including (1) DeepSeek-LLM-Chat 67B \citep{deepseek-llm}, (2) Qwen 72B \citep{qwen}, (3) SeaLLM-v2 7B \citep{seallm}, and (4) ChatGLM3 6B \citep{chatglm3}; alongside several models with particular improvements for mathematical capabilities, such as (5) InternLM2-Math 20B~\footnote{\url{https://github.com/InternLM/InternLM-Math}}, which extends InternLM2 through specialized math-centric training and subsequent instruction tuning, (6) Math-Shepherd-Mistral 7B, which introduces PPO training \citep{schulman2017proximal} to Mistral 7B \citep{mistral} by leveraging a process-supervised reward model, (7) the WizardMath suite \citep{wizardmath}, enhancing mathematical problem-solving in Mistral 7B and Llama-2 70B \citep{llama2} with evolve-instruct—a form of instruction tuning reliant on AI-generated instructions—and PPO training; this approach primarily draws training data from GSM8K and MATH, (8) MetaMath 70B \citep{metamath}, which builds upon Llama-2 70B, refined using an expanded set of GSM8K and MATH datasets, (9) ToRA 34B \cite{tora}, which adapts CodeLlama 34B for mathematical reasoning requiring integration with external tools, and (10) MAmmoTH 70B \citep{MathInstruct}, developed by further instruction tuning Llama-2 70B on the MathInstruct corpus.

Table \ref{tab:sft_rl_math} illustrates that in the absence of tool use, DeepSeekMath-Instruct 7B exhibits robust step-by-step reasoning capabilities. On the competition-level MATH dataset, our model achieves results that outperform all other open-source models and most proprietary counterparts—such as Inflection-2 and Gemini Pro—by at least 9\% in absolute terms. This performance advantage holds even when compared against models with significantly higher parameter counts (e.g., Qwen 72B) or those fine-tuned with reinforcement learning focused on mathematical problem-solving (e.g., WizardMath-v1.1 7B). Although DeepSeekMath-Instruct is competitive with leading Chinese proprietary models like GLM-4 and Baichuan-3 on the MATH benchmark, it still does not reach the performance levels of GPT-4 or Gemini Ultra.

When evaluated in a framework that permits the combination of natural language reasoning and programmatic tool application for solving problems, DeepSeekMath-Instruct 7B achieves nearly 60\% accuracy on MATH, outperforming all previously released open-source models. Across additional benchmarks, our model demonstrates performance on par with DeepSeek-LLM-Chat 67B, which represents the former state-of-the-art despite being ten times larger in scale.
\section{Reinforcement Learning}

\subsection{Group Relative Policy Optimization} After the Supervised Fine-Tuning (SFT) phase, reinforcement learning (RL) methods have demonstrated notable success in enhancing the mathematical reasoning skills of LLMs \citep{wang2023math,wizardmath}. Here, we present our approach, Group Relative Policy Optimization (GRPO), a reinforcement learning algorithm designed for both efficiency and effectiveness.

\subsubsection{From PPO to GRPO} Proximal Policy Optimization (PPO) \citep{schulman2017proximal} is an actor-critic RL algorithm that is widely used in the RL fine-tuning stage of LLMs \citep{ouyang2022training}. In particular, it optimizes LLMs by maximizing the following surrogate objective:

\begin{equation}
\footnotesize
    \mathcal{J}_{PPO}(\theta) = \mathbb{E}{[q \sim P(Q), o \sim \pi_{\theta_{old}}(O|q)]} \frac{1}{|o|} \sum_{t=1}^{|o|} \min \left[ \frac{\pi_\theta(o_{t} | q, o_{<t})}{\pi_{\theta_{old}}(o_{t} | q, o_{<t})} A_{t}, \text{clip} \left( \frac{\pi_\theta(o_{t} | q, o_{<t})}{\pi_{\theta_{old}}(o_{t} | q, o_{<t})}, 1 - \varepsilon, 1 + \varepsilon \right)  A_{t} \right] ,
\end{equation}

Here, $\pi_{\theta}$ denotes the present policy model, while $\pi_{\theta_{old}}$ refers to its predecessor. The variables $q$ and $o$ represent questions and their associated outputs, drawn both from the dataset and sampled according to the prior policy $\pi_{\theta_{old}}$. The parameter $\varepsilon$ functions as a clipping threshold, utilized in PPO to promote stable training dynamics. The term $A_t$ stands for the advantage, which is derived using Generalized Advantage Estimation (GAE) \citep{gae} and is calculated based on subsequent rewards $\{r_{\ge t}\}$ and the estimated values produced by the parameterized function $V_{\psi}$. Consequently, PPO necessitates that a value function be trained concurrently with the policy. To prevent the reward model from being excessively optimized, a prevailing method involves incorporating a per-token KL divergence penalty, computed against a reference model, into the reward at every token position \citep{ouyang2022training}, specifically,

\begin{equation}
    r_{t} = r_\varphi(q, o_{\le t}) - \beta \log\frac{\pi_{\theta}(o_{t}|q, o_{<t})}{\pi_{ref}(o_{t}|q, o_{<t})},
\label{eq:PPO-reward}
\end{equation}

Here, $r_\varphi$ denotes the reward model, while $\pi_{ref}$ represents the reference model—commonly the original SFT model. The parameter $\beta$ serves as the scaling factor for the KL divergence penalty.

As the value function employed in PPO is typically another model of comparable size as the policy model, it brings a substantial memory and computational burden. Additionally, during RL training, the value function is treated as a baseline in the calculation of the advantage for variance reduction. While in the LLM context, usually only the last token is assigned a reward score by the reward model, which may complicate the training of a value function that is accurate at each token. To address this, as shown in Figure \ref{fig:grpo}, we propose Group Relative Policy Optimization (GRPO), which obviates the need for additional value function approximation as in PPO, and instead uses the average reward of multiple sampled outputs, produced in response to the same question, as the baseline. More specifically, for each question $q$, GRPO samples a group of outputs $\{o_1, o_2, \cdots, o_G\}$  from the old policy  $\pi_{\theta_{old}}$  and then optimizes the policy model by maximizing the following objective:

\begin{equation}
\footnotesize
\begin{split}
    \mathcal{J}_{GRPO}(\theta) &= \mathbb{E}{[q \sim P(Q), \{o_i\}_{i=1}^G \sim \pi_{\theta_{old}}(O|q)]}  \\
    & \frac{1}{G}\sum_{i=1}^G\frac{1}{|o_i|} \sum_{t=1}^{|o_i|} \bigg\{ \min \left[ \frac{\pi_\theta(o_{i,t} | q, o_{i,<t})}{\pi_{\theta_{old}}(o_{i,t} | q, o_{i,<t})} \hat{A}_{i,t},\, \text{clip} \left( \frac{\pi_\theta(o_{i,t} | q, o_{i,<t})}{\pi_{\theta_{old}}(o_{i,t} | q, o_{i,<t})}, 1 - \varepsilon, 1 + \varepsilon \right) \hat{A}_{i,t} \right] \\
    & \qquad - \beta \mathbb{D}_{KL}\left[\pi_{\theta} || \pi_{ref}\right] \bigg\},
\end{split}
\label{eq:GRPO-obj}
\end{equation}

where $\varepsilon$ and $\beta$ are hyper-parameters, and $\hat{A}_{i,t}$  is the advantage calculated based on relative rewards of the outputs inside each group only, which will be detailed in the following subsections. The group relative way that GRPO leverages to calculate the advantages, aligns well with the comparative nature of rewards models,  as reward models are typically trained on datasets of comparisons between outputs on the same question. Also note that, instead of adding KL penalty in the reward, GRPO regularizes by directly adding the KL divergence between the trained policy and the reference policy to the loss, avoiding complicating the calculation of  $\hat{A}_{i,t}$. And different from the KL penalty term used in (\ref{eq:PPO-reward}), we estimate the KL divergence with the following unbiased estimator \citep{kl_approx}:

\begin{equation}
\small
    \mathbb{D}_{KL}\left[\pi_{\theta} || \pi_{ref}\right] = \frac{\pi_{ref}(o_{i,t}|q,o_{i,<t})}{\pi_{\theta}(o_{i,t}|q,o_{i,<t})}- \log\frac{\pi_{ref}(o_{i,t}|q,o_{i,<t})}{\pi_{\theta}(o_{i,t}|q,o_{i,<t})} - 1,
\end{equation}

which is guaranteed to be positive.

\begin{algorithm}[t]
  \small
  \caption{Iterative Group Relative Policy Optimization}
  \textbf{Input} initial policy model $\pi_{\theta_{\text{init}}}$; reward models $r_\varphi$; task prompts $\mathcal{D}$; 
  hyperparameters $\varepsilon$, $\beta$, $\mu$
  \begin{algorithmic}[1]
    \State Initialize policy model as $\pi_\theta \leftarrow \pi_{\theta_{\text{init}}}$
    \For{iteration = 1, \dots, I}
       \State Set reference model $\pi_{ref} \leftarrow \pi_{\theta}$
      \For{step = 1, \dots, M}
      \State Select a batch $\mathcal{D}_b$ sampled from $\mathcal{D}$
      \State Store previous policy model $\pi_{\theta_{old}} \leftarrow \pi_{\theta}$ 
      \State For every question $q \in \mathcal{D}_b$, generate $G$ responses $\{o_i\}_{i=1}^G \sim \pi_{\theta_{old}} (\cdot \mid q)$
      \State Evaluate rewards $\{r_i\}_{i=1}^{G}$ for all responses $o_i$ with $r_\varphi$ 
      \State For each $o_i$, estimate $\hat{A}_{i,t}$ for token $t$ using the group relative advantage approach
      \For{GRPO iteration = 1, \dots, $\mu$}
        \State Perform a policy update to $\pi_{\theta}$ by maximizing the GRPO objective (Equation \ref{eq:GC-GRPO})
      \EndFor
    \EndFor 
    \State Adjust $r_\varphi$ with ongoing training via a replay buffer
    \EndFor 
  \end{algorithmic}
  \textbf{Output} $\pi_\theta$
  \label{alg:iter-grpo}
\end{algorithm}

\subsubsection{Outcome Supervision RL with GRPO} Specifically, given a question $q$, a set of responses $\{o_1, o_2, \cdots, o_G\}$ is generated by the previous policy $\pi_{\theta_{old}}$. Each generated output is evaluated by a reward model, resulting in $G$ reward scores $\mathbf{r}=\{r_1, r_2, \cdots, r_G\}$. To standardize these rewards, each is centered and scaled by subtracting the mean and dividing by the standard deviation within the group. Under outcome supervision, only the final token in each response $o_i$ receives this standardized reward, and the computed advantage for all tokens in $o_i$ is uniformly set to the normalized value; formally, $\hat{A}_{i, t} = \widetilde{r}_i = \frac{r_i- {\rm mean}(\mathbf{r})}{{\rm std}(\mathbf{r})}$. The policy is then trained by maximizing the objective function given in equation (\ref{eq:GRPO-obj}).

\subsubsection{Process Supervision RL with GRPO} Outcome supervision typically delivers feedback only upon the completion of each output, which might be inadequate or inefficient for guiding policy learning in intricate mathematical scenarios. Building upon the approach in \cite{wang2023math}, we investigate process supervision as an alternative, providing evaluative signals after every individual reasoning step. More precisely, considering a question $q$ and $G$ sampled output sequences $\{o_1, o_2, \cdots, o_G\}$, a process reward model assigns step-level rewards for each output, resulting in the following reward structure: $\mathbf{R} = \{ \{r_1^{index(1)},\cdots,r_1^{index(K_1)}\}, \cdots,  \{r_G^{index(1)},\cdots,r_G^{index(K_G)}\} \}$. Here, $index(j)$ indicates the position of the concluding token of the $j$-th reasoning step, and $K_i$ denotes the number of steps in the $i$-th output. All rewards undergo normalization using their overall mean and standard deviation, given by $\widetilde{r}_i^{index(j)} = \frac{r_i^{index(j)} - {\rm mean(\mathbf{R})}}{{\rm std(\mathbf{R})}}$.
Next, for process supervision, the advantage assigned to each token is computed as the cumulative normalized rewards from all subsequent steps: $\hat{A}_{i, t} = \sum_{index(j) \ge t} \widetilde{r}_i^{index(j)}$.
Finally, the policy is updated through maximization of the objective function specified in equation (\ref{eq:GRPO-obj}).

\subsubsection{Iterative RL with GRPO} Over the course of reinforcement learning, previously trained reward models may become inadequate for effectively guiding the evolving policy model. To address this, we investigate an iterative RL framework employing GRPO. As depicted in Algorithm \ref{alg:iter-grpo}, this iterative approach involves producing new datasets for the reward model by sampling from the current policy model. The reward model is then updated through continuous training, using a replay method that retains 10\% of its prior data. Subsequently, the policy model itself becomes the reference model, and it is further trained using the updated reward model.

\subsection{Training and Evaluating DeepSeekMath-RL}

Our reinforcement learning experiments utilize DeepSeekMath-Instruct 7B as the base model. The RL training dataset is comprised exclusively of approximately 144K chain-of-thought-style questions drawn from the GSM8K and MATH portions of the SFT dataset. By omitting other SFT-derived questions, we aim to better isolate the effects of RL training on benchmarks that remain unexposed to the relevant data throughout this phase. The reward model’s training dataset is created in accordance with the methodology outlined in \citep{wang2023math}. The initial reward model is built upon DeepSeekMath-Base 7B, using a learning rate of $2\times10^{-5}$. For the GRPO method, the policy model is trained with a learning rate of $1\times10^{-6}$, and the KL regularization coefficient is set to $0.04$. Each question in the RL process is associated with $64$ sampled responses. The model is trained with a maximum sequence length of 1024 tokens and a batch size of 1024. Following every exploration interval, the policy model parameters undergo a single update step. DeepSeekMath-RL 7B is evaluated across benchmarks employing the same protocol as DeepSeekMath-Instruct 7B. Within these evaluations, GSM8K and MATH chain-of-thought problems are treated as in-domain, while performance on all remaining benchmarks is considered representative of out-of-domain generalization.

Table \ref{tab:sft_rl_math} presents a comparative analysis of open- and closed-source models employing both chain-of-thought and tool-integrated reasoning techniques on English and Chinese evaluation benchmarks. Our observations reveal the following: 1) DeepSeekMath-RL 7B achieves accuracy rates of 88.2\% on GSM8K and 51.7\% on MATH when leveraging chain-of-thought reasoning. These results outperform all open-source models ranging from 7B to 70B parameters, in addition to exceeding the accuracy of most closed-source systems. 2) Notably, DeepSeekMath-RL 7B is exclusively fine-tuned on chain-of-thought-format instruction tuning datasets from GSM8K and MATH, with DeepSeekMath-Instruct 7B serving as its initial checkpoint. Although its training data is limited in scope, DeepSeekMath-RL 7B demonstrates superior performance relative to DeepSeekMath-Instruct 7B across all evaluated metrics, thereby highlighting the benefits of reinforcement learning.

\section{Discussion}
This section presents the outcomes observed in both the pre-training and reinforcement learning experiments.

\subsection{Lessons Learnt in Pre-Training}

Initially, we present insights gained from our pre-training phase. Unless mentioned otherwise, the training configurations described in Section \ref{sec:quality-policy} are employed. Additionally, it should be highlighted that references to the DeepSeekMath Corpus within this section pertain to an 89B-token subset curated during the second round of data gathering.

\subsubsection{Code Training Benefits Mathematical Reasoning}

An often-cited but unproven theory proposes that code training enhances reasoning abilities. In this work, we seek to provide some evidence regarding this claim, focusing especially on mathematics: training on code boosts models’ mathematical reasoning skills, whether or not external tools are involved.

To study how code training affects mathematical reasoning, we experimented with the following two-stage training and one-stage training settings:

\textbf{Two-Stage Training}

\begin{itemize}[topsep=0pt]
    \item \textbf{Code Training for 400B Tokens $\rightarrow$ Math Training for 150B Tokens}: The DeepSeek-LLM 1.3B model undergoes pretraining on 400 billion code tokens, after which it is further trained on 150 billion math-specific tokens;
    \item \textbf{General Training for 400B Tokens $\rightarrow$ Math Training for 150B Tokens}: For comparison, a control setup is implemented in which, during the initial training stage, general tokens—randomly drawn from DeepSeek-AI’s extensive general corpus—replace the code tokens. This is intended to evaluate whether code tokens offer specific benefits in enhancing mathematical reasoning relative to general tokens.
\end{itemize}

\textbf{One-Stage Training}

\begin{itemize}[topsep=0pt]
    \item \textbf{Math Instruction on 150B Tokens}: DeepSeek-LLM 1.3B is trained using 150B tokens from mathematical datasets;
    \item \textbf{Combined Training on 400B Code Tokens and 150B Math Tokens}: We observe that performing math instruction after code training negatively affects coding abilities. Therefore, we explore whether training concurrently on both code and math tokens enhances mathematical reasoning while also reducing the risk of catastrophic forgetting.
\end{itemize}

\paragraph{Results}
The downstream results obtained from various training configurations are presented in Table \ref{tab:code-math} and Table \ref{tab:code-math-general-eval}.

Code training enhances program-aided mathematical reasoning across both two-stage and one-stage training paradigms. Table \ref{tab:code-math} demonstrates that, in the two-stage setting, code training alone leads to marked improvements in solving GSM8K and MATH tasks using Python. Subsequent math training during the second stage brings additional gains. Notably, in the one-stage training approach, jointly presenting code and math tokens reduces the risk of catastrophic forgetting—a problem characteristic of two-stage training—and also fosters positive interactions between coding proficiency (Table \ref{tab:code-math-general-eval}) and program-aided mathematical reasoning capabilities (Table \ref{tab:code-math}).

Training with code additionally enhances mathematical reasoning abilities even in the absence of tool use. Within the two-stage training framework, the early phase of code training yields noticeable improvements. This stage also accelerates learning during the later mathematics-focused phase, culminating in optimal results. Conversely, merging code tokens and math tokens in a single-stage training procedure diminishes performance in mathematical reasoning when tools are not involved. One possible explanation is that DeepSeek-LLM 1.3B's relatively small model size constrains its ability to absorb both code and mathematical information in parallel.

\subsubsection{ArXiv Papers Seem Ineffective in Improving Mathematical Reasoning}

ArXiv papers are commonly included as a component of math pre-training data \citep{minerva,gpt-f,llemma,mathpile}. However, detailed analysis regarding their impact on mathematical reasoning has not been extensively conducted. Perhaps counter-intuitively, according to our experiments, arXiv papers seem ineffective in improving mathematical reasoning. We experiment with models of different sizes, including DeepSeek-LLM 1.3B and DeepSeek-Coder-Base-v1.5 7B \citep{deepseek-coder}, using arXiv corpora that underwent varied processing pipelines:

\begin{itemize}[topsep=0pt]
    \item \textbf{MathPile} \citep{mathpile}: a dataset comprising 8.9 billion tokens assembled using heuristic-based data cleaning and filtering, with scientific arXiv articles constituting more than 85\% of its content;
    \item \textbf{ArXiv-RedPajama} \citep{redpajama}: a 28.0 billion token collection consisting of all arXiv LaTeX documents, excluding preambles, comments, macros, and reference sections.
\end{itemize}

For our experiments, we independently trained DeepSeek-LLM 1.3B with 150B tokens and DeepSeek-Coder-Base-v1.5 7B with 40B tokens on each individual arXiv dataset. The findings suggest that incorporating arXiv papers does not significantly enhance models' mathematical reasoning capabilities. Training solely on the arXiv corpus resulted in neither substantial progress nor, in some cases, even worsened performance on a variety of mathematical evaluation sets of varying difficulty included in this work. These evaluation sets span quantitative reasoning tasks such as GSM8K and MATH (Table \ref{tab:arxiv-cot}), multiple-choice assessments like MMLU-STEM (Table \ref{tab:arxiv-cot}), and formal mathematical reasoning benchmarks such as miniF2F (Table \ref{tab:arxiv_atp}).

However, this conclusion has its limitations and should be taken with a grain of salt. We have not yet studied:

\begin{itemize}[topsep=0pt]
    \item How arXiv tokens influence math-specific applications absent from this study, for instance, transforming formal statements or proofs into their informal counterparts (informalization of theorems);
    \item The influence of integrating arXiv tokens with additional data types;
    \item If the advantages associated with arXiv papers persist or are amplified when scaling up to larger models.
\end{itemize}

Thus, further exploration is required, which we leave for future studies.

\subsection{Insights of Reinforcement Learning}
\subsubsection{Towards to a Unified Paradigm} Here, we introduce a unified framework for examining various training strategies, including SFT, RFT, DPO, PPO, GRPO, and conduct empirical studies to investigate the components within this unified approach. In general, the gradient of a training objective with respect to the parameter $\theta$ for these methods can be formulated as follows:

\begin{equation}
    \nabla_{\theta}\mathcal{J}_{\textcolor{red}{\mathcal{A}}}(\theta) = \mathbb{E}[\underbrace{(q,o) \sim \textcolor{red}{\mathcal{D}}}_{Data \ Source}]\left( \frac{1}{|o|} \sum_{t=1}^{|o|}  \underbrace{GC_{{\mathcal{A}}}(q, o, t, \textcolor{red}{\pi_{{rf}}})}_{Gradient \ Coefficient}  \nabla_{\theta}\log \pi_{\theta}(o_t | q, o_{<t})\right).
\label{eq:objective}
\end{equation}

There exist three key components: 1) \textit{Data Source $\mathcal{D}$}, which determines the training data; 2) \textit{Reward Function $\pi_{{rf}}$}, which is the source of the training reward signal; 3) \textit{Algorithm $\mathcal{A}$}: which processes the training data and the reward signal to the gradient coefficient $GC$ that determines the magnitude of the penalty or reinforcement for the data. We analyze several representative methods based on such a unified paradigm:

\begin{itemize}
    \item  \textbf{Supervised Fine-tuning (SFT)}: SFT adapts a pretrained model to a downstream task by training it on human-curated SFT datasets.
    \item \textbf{Rejection Sampling Fine-tuning (RFT)}: RFT applies an additional stage of fine-tuning to the SFT model, using a subset of its generated responses in reply to SFT prompts. Only those responses assessed as correct are used to update the model.
    \item \textbf{Direct Preference Optimization (DPO)}: DPO enhances the SFT model by conducting further fine-tuning on augmented data, which is generated by sampling responses from the SFT model and using a pairwise DPO objective for optimization.
    \item \textbf{Online Rejection Sampling Fine-tuning (Online RFT)}: Unlike traditional RFT, Online RFT starts by initializing the policy model with the SFT weights and then incrementally updates it using augmented samples that are generated on-the-fly from the current policy model.
    \item \textbf{PPO/GRPO}: PPO/GRPO initializes the policy model from the SFT checkpoint and subsequently improves the policy via reinforcement learning, leveraging samples that are dynamically produced by the model during training.
\end{itemize}

Table \ref{tab:data_policy} presents a summary of the elements comprising these methods. For an in-depth explanation of the derivation steps, consult Appendix \ref{app:analysis-rl}.

\paragraph{Observation about Data Source} We categorize data sources as either online sampling or offline sampling. Online sampling refers to training data obtained through exploration conducted by the actively updated policy model during training. In contrast, offline sampling utilizes training data drawn from the output of the fixed, original SFT model. Methods such as RFT and DPO employ the offline approach, whereas Online RFT and GRPO utilize the online approach.

As illustrated in Figure \ref{fig:rl-analysis}, our results indicate that Online RFT achieves substantially better performance than RFT across both benchmarks. Notably, Online RFT matches RFT's effectiveness during the early phases of training, but establishes a clear lead in the subsequent stages, underscoring the benefits of the online training approach. This pattern aligns with expectations: during initial training, the actor and SFT model remain largely similar, resulting in sampled data with negligible variation. As training progresses, the actor's output starts to diverge more prominently, and the advantages of dynamically sampled data become increasingly pronounced.

\paragraph{Observation about Gradient Coefficient} The algorithm utilizes the reward signal to determine the gradient coefficient used in updating the model parameters. In our experimental setup, we decompose the reward function into two components: ‘Rule’ and ‘Model’. The ‘Rule’ component evaluates responses based solely on answer accuracy, whereas the ‘Model’ aspect involves the use of a trained reward model to assign scores to each response. This reward model is trained using labels derived from rule-based judgments. As delineated in Equations \ref{eq:GC-RFT} and \ref{eq:GC-GRPO}, there exists a notable distinction between Online RFT and GRPO: only GRPO modifies its gradient coefficient in response to the reward scores produced by the reward model. Consequently, GRPO provides nuanced reinforcement or punishment to responses relative to their associated reward magnitudes. Conversely, Online RFT does not possess this adaptability; it fails to penalize incorrect responses and instead applies identical reinforcement strength to all responses deemed correct.

As illustrated in Figure \ref{fig:rl-analysis}, GRPO outperforms online RFT, underscoring the advantages of adjusting the coefficients for positive and negative gradients. Moreover, GRPO+PS achieves better results than GRPO+OS, suggesting that employing fine-grained, step-specific gradient coefficients is beneficial. Additionally, we investigate iterative RL by performing two successive iterations in our experiments. Figure \ref{fig:iter-rl} demonstrates that iterative RL leads to marked performance gains, with the most pronounced improvement occurring during the initial iteration.

\subsubsection{Why RL Works?} In this work, we apply reinforcement learning to a selected portion of instruction tuning datasets and observe considerable performance gains over models trained only with instruction tuning. To elucidate the underlying reasons for reinforcement learning's effectiveness, we analyze the Pass@K and Maj@K accuracies for both Instruct and RL-enhanced models across two evaluation benchmarks. As depicted in Figure \ref{fig:maj-pass}, reinforcement learning leads to improvements in Maj@K, but does not substantially affect Pass@K scores. This suggests that reinforcement learning primarily contributes to the overall robustness of the output distribution, that is, \textbf{the observed improvements stem from elevating the most accurate response within the TopK outputs rather than by improving the model's inherent competence}. Likewise, \citep{wang2023making} have pointed out a \textbf{misalignment problem} in reasoning tasks encountered in SFT models and demonstrated that a variety of preference alignment approaches \citep{yuan2023rrhf,song2023preference,wang2023making} can yield enhanced reasoning abilities in such models.

\subsubsection{How to Achieve More Effective RL?}
\label{sec:effec_rl}
Our results show that RL performs admirably on mathematical reasoning tasks. Furthermore, we introduce a cohesive framework that clarifies how various notable training approaches fit within the RL landscape. This framework enables us to view these methods as manifestations or approximations of RL principles. As outlined in Equation \ref{eq:objective}, three primary elements underpin these methods: the Data Source, the Algorithm, and the Reward Function. In what follows, we highlight several promising research directions related to each of these key elements.

\paragraph{Data Source} The choice of data source serves as the foundational input for any training methodology. Within RL settings, we define the data source as unlabeled questions paired with responses generated from the policy model. For this study, we restrict ourselves to utilizing questions derived from the instruction tuning phase and employ straightforward nucleus sampling to obtain model outputs. We hypothesize that this methodological constraint might underlie the observation that our RL pipeline yields improvements only in Maj@K metrics. For subsequent work, we plan to assess the RL pipeline using out-of-distribution prompts and integrate \textbf{advanced sampling (decoding) strategies}, such as those employing tree-search techniques \citep{yao2023tree}. Furthermore, \textbf{efficient inference techniques} \citep{xia-etal-2023-speculative,leviathan2023fast,kwon2023efficient,xia2024unlocking}, which shape how effectively the policy model explores the space of possible outputs, represent an additional pivotal factor.

\paragraph{Algorithms} Algorithms use data along with the reward signals to determine the direction and magnitude of gradient updates for the model parameters. According to Equation \ref{eq:objective}, contemporary approaches generally place complete \textbf{TRUST} in the reward function's signals to modulate the conditional likelihood of specific tokens, either increasing or decreasing them as needed. Nonetheless, it remains unfeasible to guarantee the consistent reliability of reward signals, particularly when addressing highly complex tasks. For instance, the PRM800K datasets \citep{lightman2023let}, which underwent meticulous annotation by skilled annotators, still contain roughly 20\% erroneous annotations\footnote{\url{https://github.com/openai/prm800k/issues/12\#issuecomment-1728491852}}. Therefore, our focus is on investigating reinforcement learning algorithms capable of maintaining robustness in the presence of noisy reward information. We contend that such \textbf{WEAK-TO-STRONG} \citep{burns2023weak} alignment strategies are poised to fundamentally transform learning paradigms.

\paragraph{Reward Function} The reward function serves as the foundational training signal in RL and is typically instantiated as a neural reward model. We identify three key research avenues concerning reward models: 1) \textbf{Improving the generalization capability of the reward model.} It is essential for the reward model to generalize robustly to unseen questions and complex outputs generated during advanced decoding; if it fails to do so, reinforcement learning may simply preserve the status quo of LLM distributions instead of genuinely enhancing their abilities; 2) \textbf{Capturing and leveraging uncertainty in reward models.} Characterizing uncertainty may help bridge the gap between underpowered reward models and the learning algorithms that transition from weaker to stronger performance; 3) \textbf{Developing efficient methods for constructing high-quality process reward models,} which can deliver granular training feedback for tasks that require stepwise reasoning \citep{lightman2023let,wang2023math}.

\section{Conclusion, Limitation, and Future Work} 
We introduce DeepSeekMath, an open-source model that surpasses all other open-source competitors on the MATH benchmark at the competition level and comes close to matching the effectiveness of proprietary systems. DeepSeekMath is built upon DeepSeek-Coder-v1.5 7B as its foundation and is further trained on a corpus totaling 500B tokens, prominently including 120B math-specific tokens extracted from Common Crawl resources. Our detailed ablation analyses indicate that web-based materials are a substantial reservoir of high-quality mathematical content, whereas the anticipated value from arXiv content is less pronounced. We propose Group Relative Policy Optimization (GRPO), an adaptation of Proximal Policy Optimization (PPO), demonstrating that GRPO significantly enhances mathematical reasoning skills and reduces memory usage. Experimental evaluations confirm GRPO’s utility even when DeepSeekMath-Instruct 7B attains high benchmark results. Additionally, we present a cohesive framework for interpreting various related approaches and outline prospective avenues for advancing reinforcement learning methodologies.

While DeepSeekMath~demonstrates strong results on benchmarks involving quantitative reasoning, its performance in geometry-related tasks and theorem-proving remains inferior to that of proprietary models. For example, during preliminary experiments, the model was unable to solve questions concerning triangles and ellipses, which could be symptomatic of a bias in data selection during both pre-training and fine-tuning stages. Furthermore, the limited model size constrains DeepSeekMath~'s performance in few-shot scenarios: unlike GPT-4, whose results improve when presented with few-shot examples, DeepSeekMath~yields similar outcomes under both zero-shot and few-shot conditions. Looking ahead, we plan to enhance our data engineering procedures to assemble a pre-training corpus of higher quality. Additionally, we will investigate avenues for more effective reinforcement learning methodologies for LLMs, as outlined in Section \ref{sec:effec_rl}.

\end{document}